% Toolkits compared with aeon: engines, wrappers and defaults.
% Standalone; to use in the paper, copy the section and table body.
\documentclass[11pt]{article}
\usepackage[margin=2.5cm]{geometry}
\usepackage{booktabs}
\usepackage{tabularx}
\usepackage{array}

\begin{document}

\section*{Toolkits compared}

We compare aeon with six other Python toolkits, which between them contain
five distinct estimation engines (Table~\ref{tab:toolkits}). statsforecast and
skforecast reimplement the algorithms of R's \texttt{forecast} package,
statsforecast with a compiled C++ core and skforecast in numba. smooth is a
Python port of the ADAM framework of R's \texttt{smooth} package, also with a
C++ core. statsmodels provides state-space ARIMA and ETS models estimated by
Kalman-filter maximum likelihood, but no automatic model selection, so we
supply a grid search by AIC that mirrors aeon's candidate set. The remaining
two toolkits are wrappers. pmdarima implements the stepwise search of R's
\texttt{auto.arima} but fits every candidate with statsmodels'
\texttt{SARIMAX}. sktime's \texttt{AutoARIMA} wraps pmdarima, and its
\texttt{AutoETS} runs its own selection loop over statsmodels'
\texttt{ETSModel}. Given the true model, sktime therefore produces the same
forecasts as pmdarima for ARIMA and as statsmodels for ETS. aeon's
implementations are written from scratch in numba: ARIMA is fitted by
conditional sum of squares and ETS by maximum likelihood, both optimised by
Nelder--Mead with multiple starts, with models selected by AIC.

The toolkits differ in defaults as well as engines, and we kept each toolkit's
defaults except where a common setting was needed for a fair comparison. Every
toolkit was given the true seasonal period and the same auto-ARIMA order
limits, and seasonal ARIMA was disabled throughout because aeon's ARIMA is
non-seasonal. The remaining differences are part of what a user of each
toolkit gets. statsforecast and smooth select models by AICc, the others by
AIC. smooth initialises its states by backcasting, while the others optimise
them. Following R, skforecast fits the candidates in its ARIMA search by
conditional sum of squares when the series has more than 150 observations, and
refits only the selected model by maximum likelihood. This is a likely reason
why it overtakes aeon on long series while remaining equally accurate: aeon
also searches with conditional sum of squares, but fits each candidate from
several Nelder--Mead starts, with and without a constant, where skforecast
runs a single gradient-based fit.

\begin{table}[ht]
\centering
\small
\caption{Automatic ARIMA and ETS in the toolkits compared, and the estimation
engine each relies on. Versions: aeon 8585c7bf, statsforecast 2.1.1,
skforecast 0.26.0, pmdarima 2.1.1, sktime 1.2.0, statsmodels 0.15.0,
smooth 1.2.0.}
\label{tab:toolkits}
\begin{tabularx}{\textwidth}{@{}l>{\raggedright\arraybackslash}p{2.6cm}>{\raggedright\arraybackslash}p{2.6cm}>{\raggedright\arraybackslash}X@{}}
\toprule
Toolkit & Auto ARIMA & Auto ETS & Engine \\
\midrule
aeon & own & own &
  numba; ARIMA by conditional sum of squares, ETS by likelihood, both with
  multi-start Nelder--Mead; AIC \\
statsforecast & own & own &
  port of R \texttt{forecast} (\texttt{auto.arima}, \texttt{ets}) with a
  compiled C++ core; AICc \\
skforecast & own (\texttt{Arima}) & own (\texttt{Ets}) &
  numba reimplementation following R \texttt{forecast}; CSS search for
  series longer than 150 \\
pmdarima & own search & -- &
  \texttt{auto.arima}-style stepwise search; each fit is statsmodels
  \texttt{SARIMAX} (Kalman filter) \\
sktime & wraps pmdarima & own loop over statsmodels \texttt{ETSModel} &
  statsmodels \\
statsmodels & none (grid search supplied) & none (grid search supplied) &
  Kalman-filter \texttt{ARIMA} and \texttt{ETSModel}, gradient-based maximum
  likelihood \\
smooth & \texttt{AutoMSARIMA} & \texttt{ES} &
  port of R \texttt{smooth} (ADAM); C++ core, \texttt{nlopt} optimiser,
  AICc, initial states by backcasting \\
\bottomrule
\end{tabularx}
\end{table}

\end{document}
